\centerline{\bf \Huge Пересчитаем ребра II}

\begin{flushright}
        \scriptsize{}
\end{flushright}

\problem{1}
У каждого из 16   детей было по 8   конфет. Каждую минуту один из детей платит в кассу столько рублей, у скольких детей конфет не меньше, чем у него, а затем съедает одну свою конфету. Сколько денег может оказаться в кассе, когда все конфеты будут съедены?


\problem{2}
В углах доски $3 \times 3$ стоят кони: по 2 коня черного и белого цветов, причем кони одного цвета стоят в противоположных углах. За один ход разрешается выбрать любого коня и сделать им ход в свободную клетку. Можно ли переставить коней так, чтобы по прежнему все кони стояли в углах, но кони одного цвета стояли в углах при одной стороне?

\problem{3}
В королевстве живут графы, герцоги и маркизы. Однажды каждый граф сразился на дуэли с тремя герцогами и несколькими маркизами. Каждый герцог сразился на дуэли с двумя графами и шестью маркизами. Каждый маркиз сразился на дуэли с тремя герцогами и двумя графами. Известно, что все графы сразились с равным числом маркизов. Со сколькими маркизами сразился каждый граф?

\problem{4}
В классе больше 30, но меньше 35 человек. Любой мальчик дружит с тремя девочками, а любая девочка – с пятью мальчиками. Сколько человек в классе?

\problem{5}
В классе 20 учеников, причём каждый дружит не менее, чем с 14 другими.
Можно ли утверждать, что найдутся четыре ученика, которые все дружат между собой?

\problem{6}
Есть две страны: Обычная и Зазеркалье. У каждого города в Обычной стране есть двойник в Зазеркалье и наоборот. Однако, если в Обычной стране какие-то два города соединены авиалинией, то в Зазеркалье эти города не соединены, а два любых несоединенных в Обычной стране города обязательно соединены авиалинией в Зазеркалье. В Обычной стране Алиса не может добраться из Бристоля в Ливерпуль, сделав менее двух пересадок. Докажите, что в Зазеркалье Алиса сможет перелететь из любого города в любой другой, сделав не более двух пересадок.
